\begin{abstract}
Vision-language models (VLMs) are increasingly deployed in safety-critical settings, yet their safety alignment is still evaluated almost exclusively on clean images. Existing robustness studies primarily examine general utility under common real-world corruptions, while the limited safety-focused work often considers benchmarks in which harmful content is embedded as rendered text rather than conveyed through the semantic content of the image. Consequently, it remains unclear whether VLM safety mechanisms remain reliable when safety-relevant visual information is affected by blur, noise, compression, lighting changes, and other common image degradations. In this work, we present a systematic evaluation of the safety robustness of both reasoning and non-reasoning VLMs across a diverse set of image corruptions and severity levels. We evaluate standard and safety-aligned models on multimodal safety benchmarks in which harmful intent is conveyed through image semantics, or the joint interpretation of the image and prompt. %Crucially, we distinguish unsafe compliance, safety-based refusal, and perceptual failure, while separately measuring whether the model still understands the safety-relevant visual content.
Our results reveal a safety–capability decoupling: mild corruptions systematically increase attack success rates even when performance on standard utility benchmarks is fully preserved --- corruptions with zero measured utility cost still weaken safety. The direction of the effect follows a simple harm-location principle: corruption degrades safety when the image carries the evidence of harm, and improves it when the image carries the attack itself. Existing safety alignment remains vulnerable --- corruption re-opens part of the gap that safety tuning closes --- and a module ablation localizes the failure: the refusal decision resides in the language model, so corruption-aware training of the vision tower alone yields safer reasoning but not safer answers. Finally, we evaluate prompting as a lightweight mitigation across three models: a generic safety prompt recovers more safety than corruption removes, at no measurable cost on ScienceQA or MathVista, but the corruption-aware variant that works best on the model it was developed on is worse than the generic prompt on both held-out models --- a defense that a single-model evaluation would have reported as general. These findings identify common image corruptions as an overlooked failure mode in multimodal safety evaluation and motivate corruption-aware safety benchmarking for real-world VLM deployment.

\end{abstract}
